\documentclass{article}
\usepackage{graphicx} % Required for inserting images

\usepackage{amsthm}

\newtheorem{definition}{Definition}
\newtheorem{proposition}{Proposition}

\title{Entropy}
\author{Thanh Vuong Dang}
\date{25/05/2026}

\begin{document}

\maketitle

Following the well-known paper of C.E.Shannon $\cite{shannon}$, we present some fundamental properties of entropy in information theory.
\begin{definition}
Given a set of possible events with probability occurrence $p_{1}, p_{2}, ... p_{n}.$ The entropy $H$ is defined as
\begin{equation}
H = -K \sum_{i= 1}^{n}p_{i}\log(p_{i}), 
\end{equation}
where $K$ is a positive constant. \\
Denote $H = H(p_{1}, p_{2}, ..., p_{n})$
\end{definition}
With the formula in the definition, we have some properties regarding to entropy (mentioning in Wikipedia \cite{wikipedia_entropy_information_theory}).  
\begin{proposition}
$H$ is satisfied some properties. 
\begin{itemize}
    \item \textbf{Continuity:} $H$ is continuous with respect to all $p_{i}$. 
    \item \textbf{Symmetry:} The value of $H$ is unchanged for any permutations of $p_{i}$.
    \item \textbf{Maximum:} $H$ is maximum when all probabilities $p_{i}$ are equal. 
    \item \textbf{Increasing property:} Under the condition of equiprobable events, the entropy increase when all outcomes increase, i.e  
    $$H\left(\frac{1}{n}, \frac{1}{n}, ..., \frac{1}{n}\right) < H\left(\frac{1}{n+1}, \frac{1}{n+1}, ..., \frac{1}{n+1}\right).$$
    \item \textbf{Additivity:}
\end{itemize}
\end{proposition}
The properties of entropy for random variables $X, Y$ is described as
\begin{proposition}
\begin{itemize}
    \item \textbf{Subadditivity:} $H(X, Y) \leq H(X) + H(Y)$ for jointly distributed random variables $X, Y$. 
    \item \textbf{Additivity:} $H(X, Y) = H($.
    \item \textbf{Expansibility:} $H$ is maximum when all probabilities $p_{i}$ are equal. 
    \item \textbf{Symmetry:} Under the condition of equiprobable events, the entropy increase when all outcomes increase, i.e  
    $$H\left(\frac{1}{n}, \frac{1}{n}, ..., \frac{1}{n}\right) < H\left(\frac{1}{n+1}, \frac{1}{n+1}, ..., \frac{1}{n+1}\right).$$
    \item \textbf{Small for small probability:} $\lim_{q \to 0}H(q,1-q) = 0.$
\end{itemize}
\end{proposition}

\bibliographystyle{plain}
\bibliography{references}
\end{document}
